\section{Stopping the propagation of normalized defects}\label{stop:section}\label{rec:martingale}

We return to the root-sum norm $N_\Sigma$ on the infinite tree from
Section~\ref{rec:norms}.  This section gives a second proof that a tail
of fixed size forces a uniform midpoint gain.  Its main object is a
product of ratios along a branch.  We stop the product when either a
normalized base field becomes small, the product becomes small, or an endpoint
defect becomes large.  A finite martingale controls the second kind of
stop; the other two are controlled directly by the fields.

\begin{lemma}[Uniform sequential separation]\label{stop:sequential}
For each $n$, let $x_n,z_n$ be finitely supported real vectors on
$\mathbb N^{<\omega}$, with $N_\Sigma(x_n)=1$.  Let $H_n$ be the
predecessor closure of $\operatorname{supp}x_n$, and suppose that
$(z_n)_s=0$ for every $s\in H_n$.  Then
\[
 \frac{N_\Sigma(x_n+z_n)+N_\Sigma(x_n-z_n)}2\longrightarrow1
 \quad\Longrightarrow\quad N_\Sigma(z_n)\longrightarrow0.
\]
In particular, the assertion is uniform in the height and size of $H_n$.
\end{lemma}

\begin{proof}
We first work with one pair $x,z$, writing $H$ for its head and omitting
$n$.  Write $C_s^1=P_s(x)$ and $C_s^2=Q_s(x)$.  For $s\in H$ define
\begin{equation}\label{stop:normalization}
 w_s=(C_s^1)^2+(C_s^2)^2,\qquad
 \nu_s=\frac{w_s}{w_o},\qquad
 u_s^i=\frac{C_s^i}{\sqrt{w_s}},\qquad
 \ell_s=\min\{u_s^1,u_s^2\}.
\end{equation}
Every subtree rooted in $H$ contains a nonzero coordinate of $x$.
The child inequalities therefore imply $w_s>0$.  At the root,
$1/2\le w_o\le1$, since $P_o(x)+Q_o(x)=1$.
Also $(u_s^1)^2+(u_s^2)^2=1$.  Let $I_s$ be the children of $s$ in
$H$, and $J_s$ the children outside $H$.  Choose an index $i_s\in\{1,2\}$
for which the child inequality of the base vector is tight at $s$.

For $\sigma\in\{-1,1\}$ set
\begin{equation}\label{stop:defects}
 D_s^\sigma=P_s(x+\sigma z)-P_s(x)
             =Q_s(x+\sigma z)-Q_s(x),\qquad
 d_s^\sigma=\frac{D_s^\sigma}{\sqrt{w_s}}\quad(s\in H).
\end{equation}
The two differences agree because $z_s=0$ on $H$.  They are nonnegative:
at exterior children the base fields are zero, and upward induction
uses the monotonicity of the recursion in both child bounds.  Thus
\begin{equation}\label{stop:root-defect}
 N_\Sigma(x+\sigma z)=1+2D_o^\sigma.
\end{equation}
Here and below, $C^1(v)=P(v)$ and $C^2(v)=Q(v)$ when an argument is
specified.  Define the exterior energies by
\begin{equation}\label{stop:exterior}
 E_s^\sigma=\frac1{w_o}\sum_{c\in J_s}C_c^{i_s}(\sigma z)^2,
 \qquad
 \mathcal E=\sum_{s\in H}\sum_\sigma E_s^\sigma
 =\frac1{w_o}\sum_{s\in H}\sum_{c\in J_s}
       \bigl(P_c(z)^2+Q_c(z)^2\bigr).
\end{equation}
Negation interchanges the fields, so the last sum is independent of
our choices of $i_s$.  Only finitely many exterior terms are nonzero.
On $H$ the fields of $z$ agree.  The recursion there gives
\[
 P_s(z)^2=\sum_{c\in I_s}P_c(z)^2+
 \max\left\{\sum_{c\in J_s}P_c(z)^2,
             \sum_{c\in J_s}Q_c(z)^2\right\}.
\]
Telescoping over $H$ proves
\begin{equation}\label{stop:tail-energy}
 N_\Sigma(z)^2\le4w_o\mathcal E.
\end{equation}
We will show that a small root defect forces this exterior energy to
be small, independently of the head.

There is a second useful estimate, for every $a\in H$:
\begin{equation}\label{stop:subtree}
 \sum_{\substack{s\in H\\a\preceq s}}\sum_\sigma E_s^\sigma
 \le\nu_a\sum_\sigma(\ell_a+d_a^\sigma)^2.
\end{equation}
Choose one index $i$ with $u_a^i=\ell_a$, and iterate the child inequality
for the single field $C^i(x+\sigma z)$ from $a$ to all first exterior
nodes below it.  At those nodes the whole base subtree vanishes.
Summing over signs makes the resulting boundary sum independent of
$i$, and division by $w_o$ gives \eqref{stop:subtree}.

\medskip
\noindent\emph{The identity that propagates the defect.}
For a fixed sign abbreviate $d_s=d_s^\sigma$.  Subtracting the tight
base inequality for $C^{i_s}$ from its endpoint inequality yields
\begin{equation}\label{stop:increment}
 \nu_s(2u_s^{i_s}d_s+d_s^2)
 \ge\sum_{c\in I_s}\nu_c(2u_c^{i_s}d_c+d_c^2)+E_s^\sigma.
\end{equation}
Whenever $u_s^{i_s}>0$, put $R_{sc}=u_c^{i_s}/u_s^{i_s}$.
Tightness gives $\sum_{c\in I_s}\nu_cR_{sc}^2=\nu_s$.
Expanding the squares in the following display and using that identity
rewrites \eqref{stop:increment} as
\begin{equation}\label{stop:square}
 2(u_s^{i_s}+d_s)
 \left(\nu_sd_s-\sum_{c\in I_s}\nu_cR_{sc}d_c\right)
 \ge E_s^\sigma+
       \sum_{c\in I_s}\nu_c(d_c-R_{sc}d_s)^2.
\end{equation}
The last sum measures the failure of the defect to propagate by the
ratios $R_{sc}$.  Multiplying those ratios along a branch makes the
left side telescope.

Fix $0<\lambda<1/2$, $0<\zeta<1$, and $h>0$.  Inspect the root with
$Z_o=1$.  At an inspected node $s$, stop when
\begin{equation}\label{stop:rule}
 \ell_s<\lambda,\qquad Z_s<\zeta,
 \qquad\hbox{or}\qquad\max_\sigma d_s^\sigma>h.
\end{equation}
If none holds, call $s$ a continuation node and inspect its children
$c\in I_s$ with $Z_c=Z_sR_{sc}$.  The denominator in $R_{sc}$ is at
least $\lambda$ at such a node.  Let $G$ and $S$ denote the continuation
and stop sets.  Every leaf of $H$ has $\ell_s=0$, so every inspected
branch stops.  The stops are incomparable, and the child inequalities
give
\begin{equation}\label{stop:mass}
 \sum_{T\in S}\nu_T\le1.
\end{equation}
Multiplying \eqref{stop:square} at $s\in G$ by
$Z_s/[2(u_s^{i_s}+d_s^\sigma)]$, which is at least
$\zeta/[2(1+h)]$, and summing gives
\begin{equation}\label{stop:telescope}
 \begin{split}
 \sum_{T\in S}\nu_TZ_Td_T^\sigma
 &+\frac{\zeta}{2(1+h)}\sum_{s\in G}
 \left(E_s^\sigma+\sum_{c\in I_s}
        \nu_c(d_c^\sigma-R_{sc}d_s^\sigma)^2\right)
 \le d_o^\sigma.
 \end{split}
\end{equation}
If the root itself stops, this is equality with $G$ empty.

A nonroot stop $T$ has a parent $s\in G$, and
$R_{sT}d_s^\sigma\le h/\lambda$.  Use
$(a+b)^2\le2a^2+2b^2$, then \eqref{stop:mass} and
\eqref{stop:telescope}, to obtain
\begin{equation}\label{stop:squares}
 \sum_{T\in S}\nu_T(d_T^\sigma)^2
 \le(d_o^\sigma)^2+\frac{4(1+h)}{\zeta}d_o^\sigma
                         +2(h/\lambda)^2.
\end{equation}
The first term includes the possible root stop.  Classify stops first
by $\ell_T<\lambda$, then by $\ell_T\ge\lambda$ and $Z_T<\zeta$.
At every remaining stop, $Z_T\ge\zeta$ and some $d_T^\sigma>h$.
The boundary term in \eqref{stop:telescope} therefore gives
\begin{equation}\label{stop:minima}
 \sum_{T\in S}\nu_T\ell_T^2
 \le\lambda^2+
 \sum_{\substack{T\in S\\\ell_T\ge\lambda,\ Z_T<\zeta}}\nu_T
 +\frac{\sum_\sigma d_o^\sigma}{\zeta h}.
\end{equation}
Only the middle term still needs a bound uniform in the tree.

\medskip
\noindent\emph{The probability of a small product.}
Run the following finite random process from $o$.  At $s\in H$ move to
$c\in I_s$ with probability $w_c/w_s$, and move to absorbing states
labelled $1$ and $0$ with respective probabilities
\[
 \frac{P_s(x)^2-\sum_{c\in I_s}P_c(x)^2}{w_s},\qquad
 \frac{Q_s(x)^2-\sum_{c\in I_s}Q_c(x)^2}{w_s}.
\]
These are nonnegative and sum with the child probabilities to one.
Stay at an absorbing state and run to one step beyond the height of
$H$.  Reaching a particular node $s$ has probability $\nu_s$.
Give that node value $p_s=P_s(x)^2/w_s$, and an absorbing state its
label.  The conditional expectation of the next value is $p_s$, since
\[
 \sum_{c\in I_s}\frac{w_c}{w_s}\frac{P_c(x)^2}{w_c}
 +\frac{P_s(x)^2-\sum_{c\in I_s}P_c(x)^2}{w_s}=p_s.
\]
The successive values $p_j$ thus form a martingale in $[0,1]$ for the
filtration generated by the path.

On $[\lambda^2,1-\lambda^2]$ define
$H_1(p)=\tfrac12\log p$ and $H_2(p)=\tfrac12\log(1-p)$.
Their derivatives satisfy
\begin{equation}\label{stop:derivative-constants}
 |H_i'|\le B_\lambda:=\frac1{2\lambda^2},\qquad
 |H_i''|\le2K_\lambda:=\frac1{2\lambda^4}.
\end{equation}
At a step starting at $s\in G$ take $b_j=H_{i_s}'(p_s)$; at every
other step take $b_j=0$.  These coefficients are known before the step.
With $\Delta p_j=p_{j+1}-p_j$, put
\[
 V=\sum_j(\Delta p_j)^2,\qquad M=\sum_j b_j\Delta p_j.
\]
The sum $M$ is a predictable martingale transform~\cite{Burkholder1966}.
Here its second moment follows directly: conditional mean-zero increments
make the cross terms vanish, giving
\begin{equation}\label{stop:moments}
 \mathbb EV=\mathbb E p_{\rm end}^2-p_o^2\le1,
 \qquad
 \mathbb EM^2=\sum_j\mathbb E\bigl[b_j^2(\Delta p_j)^2\bigr]
 \le B_\lambda^2.
\end{equation}
Suppose that the path reaches $T\in S$ with $\ell_T\ge\lambda$ and
$Z_T<\zeta$.  All preceding nodes are in $G$, and all values through
$T$ lie in $[\lambda^2,1-\lambda^2]$.  On each such edge $s\to c$,
Taylor's inequality gives
\[
 \log R_{sc}=H_{i_s}(p_c)-H_{i_s}(p_s)
 \ge b_j\Delta p_j-K_\lambda(\Delta p_j)^2.
\]
After this stop the path meets no further node of $G$.  Thus the later
coefficients of $M$ are zero, and allowing later nonnegative terms in
$V$ gives $\log\zeta>\log Z_T\ge M-K_\lambda V$.
For $L=\log(1/\zeta)$ this event is contained in
$\{M<-L/2\}\cup\{K_\lambda V>L/2\}$.  Chebyshev's and Markov's
inequalities, followed by \eqref{stop:moments}, show that
\begin{equation}\label{stop:probability}
 \sum_{\substack{T\in S\\\ell_T\ge\lambda,\ Z_T<\zeta}}\nu_T
 \le r_\lambda(\zeta):=
 \frac{4B_\lambda^2}{L^2}+\frac{2K_\lambda}{L}.
\end{equation}
Different stops correspond to disjoint path events.  For each fixed
$\lambda$, the right side tends to zero as $\zeta\downarrow0$.

\medskip
\noindent\emph{Finishing the estimate in the required order.}
The head consists of $G$ and the disjoint subtrees beginning at its
stops.  By \eqref{stop:subtree}, the latter contribute at most
\[
 4\sum_{T\in S}\nu_T\ell_T^2+
 2\sum_\sigma\sum_{T\in S}\nu_T(d_T^\sigma)^2
\]
to $\mathcal E$.  Write $\alpha=\sum_\sigma d_o^\sigma$ and
$\beta_0=\sum_\sigma(d_o^\sigma)^2$.  Combining
\eqref{stop:telescope}--\eqref{stop:minima} with
\eqref{stop:probability} proves the finite estimate
\begin{equation}\label{stop:uniform-energy}
 \begin{split}
 \mathcal E\le{}&4\lambda^2+4r_\lambda(\zeta)
                  +8(h/\lambda)^2+2\beta_0\\
 &+\left(\frac{10(1+h)}{\zeta}+\frac4{\zeta h}\right)\alpha.
 \end{split}
\end{equation}
Restore the sequence index.  The assumed convergence and
\eqref{stop:root-defect}, together with $w_o\ge1/2$, imply
$d_o^\sigma\to0$ for both signs.  Keep $\lambda,\zeta,h$ fixed
while taking this limit in \eqref{stop:uniform-energy}.  Then
\[
 \limsup_n\mathcal E_n
 \le4\lambda^2+4r_\lambda(\zeta)+8(h/\lambda)^2.
\]
Choose $\lambda$ small, then choose $\zeta$ small with that
$\lambda$ fixed, and finally choose $h$ small.  The right side becomes
arbitrarily small.  Therefore $\mathcal E_n\to0$, and
\eqref{stop:tail-energy} proves the lemma.
\end{proof}

\begin{corollary}\label{stop:all-radii}
For every $t>0$, the averaged asymptotic modulus of $N_\Sigma$ is positive.
\end{corollary}

\begin{proof}
Fix $t>0$.  Lemma~\ref{stop:sequential} gives an $\eta>0$, independent
of the finite unit center $x$, such that the average endpoint norm is
at least $1+\eta$ for every finite tail $z$ vanishing on its predecessor
closure and satisfying $N_\Sigma(z)\ge t/2$.  Otherwise a sequence
would contradict the lemma; convexity always bounds the average below
by one.

Given a unit center $x_0$ in the completion, choose a finite unit vector
$x$ with $N_\Sigma(x-x_0)<\eta/2$, and let $H$ be its predecessor
closure.  The finite tails are dense in $\ker\Pi_H$, as established in
Section~\ref{rec:norms}.  A tail of norm at least $t$ is the limit of
finite tails whose norms are eventually at least $t/2$.  Continuity
therefore gives the endpoint estimate at $x$ for this completed tail.
Replacing $x$ by $x_0$ loses less than $\eta/2$ in the average.
Taking $z=ty$ with $y\in\ker\Pi_H$ and $N_\Sigma(y)\ge1$, then
performing the defining infimum and supremum, gives
$\widehat\delta_{N_\Sigma}(t)\ge\eta/2>0$.
\end{proof}

\section{Logarithmic propagation on the joined tree}\label{log:section}

We now use the norm $N_J$ on $X_J$ from
construction~\ref{rec:joined-sum-definition}: the finite-height trees are joined
at a root whose coordinate is free, and the norm is the sum of its two
fields.  The following argument begins with a finite-tree statement
about arbitrary nonnegative fields.  It propagates a positive amount
of boundary mass through a product of ratios.  Taking logarithms
turns that product into increments of two bounded supermartingales.

Call a nonnegative field $a$ on a finite rooted tree an \emph{upper
field} if $a_s^2\ge\sum_{c\in\operatorname{ch}(s)}a_c^2$ at every
nonleaf.

\begin{lemma}[Uniform logarithmic propagation]\label{log:finite-tree}
Fix $M\ge\beta>0$.  Let $p,q,r$ be nonnegative fields on a finite
rooted tree with root $o$.  Suppose that $p,q,p+r,q+r$ are upper
fields, and at every nonleaf the upper inequality for at least one of
$p,q$ is an equality.  If
\[
 p_o+q_o+2r_o\le M,\qquad \sum_{b\text{ a leaf}}r_b^2\ge\beta^2,
\]
then $r_o\ge\gamma(M,\beta)>0$, independently of the tree.
One possible choice is
\begin{equation}\label{log:gamma}
 \begin{gathered}
 h=\frac{\beta^2}{2M^2},\qquad c_* =\frac{\beta\sqrt h}{8},\qquad
 H_*=1+\log(M/c_*),\qquad
 K=\frac{4(2+\sqrt6)H_*}{h},\\
 \gamma(M,\beta)=\frac{\beta h}{8}e^{-K}.
 \end{gathered}
\end{equation}
\end{lemma}

\begin{proof}
Put $u=p+r/2$ and $v=q+r/2$.  These are upper fields by convexity.
We first distribute probability along the tree.  Set $\mu_o=1$, and
choose masses recursively so that the masses of the children of every
nonleaf sum to its mass and
\begin{equation}\label{log:masses}
 \mu_s\ge M^{-2}\bigl((p_s+r_s)^2+(q_s+r_s)^2\bigr).
\end{equation}
At the root the right side is at most one.  At every other step the
sum of the children's required lower bounds is at most the required
lower bound at their parent.  Assign those lower bounds first and
distribute the remaining mass among the children.  This constructs
the masses.  Nodes of zero mass have all three fields zero and can
be omitted from the random process.

Start a path at the root, moving from $s$ to $c$ with probability
$\mu_c/\mu_s$.  Once it reaches a leaf, leave it there until the common
terminal time $l$, the height of the tree.  Relative to the filtration
of the path, define
\[
 U_i=\frac{u_{s_i}}{\sqrt{\mu_{s_i}}},\qquad
 V_i=\frac{v_{s_i}}{\sqrt{\mu_{s_i}}},\qquad
 R_i=\frac{r_{s_i}}{\sqrt{\mu_{s_i}}}.
\]
All three processes lie in $[0,M]$, and $U_i,V_i\ge R_i/2$.
The upper inequalities give
$\mathbb E_i U_{i+1}^2\le U_i^2$ and
$\mathbb E_i V_{i+1}^2\le V_i^2$.
Conditional Cauchy--Schwarz also makes $U,V$ supermartingales.

At each nonleaf choose $U$ if the inequality for $p$ is tight, and
choose $V$ otherwise.  Write $W_i^{(i)},W_{i+1}^{(i)}$ for the two
consecutive values of the process chosen at time $i$; choose $U$ once
the path has stopped at a leaf.  Subtracting the tight squared
inequality for $p$ from the one for $p+r$ gives
\[
 u_sr_s\ge\sum_c u_cr_c.
\]
The corresponding formula holds for $v$ when $q$ is tight.
Consequently
\[
 W_i^{(i)}R_i\ge
 \mathbb E_i\bigl[W_{i+1}^{(i)}R_{i+1}\bigr].
\]
Put $L_i=W_{i+1}^{(i)}/W_i^{(i)}$ when the denominator is positive,
and $L_i=0$ otherwise.  A zero denominator forces $R_i=0$ and the
conditional product on the right to vanish, so in every case
$R_i\ge\mathbb E_i[L_iR_{i+1}]$.  Iterating yields
\begin{equation}\label{log:product}
 r_o\ge\mathbb E\left[R_l\prod_{i<l}L_i\right].
\end{equation}

The boundary hypothesis says
$\mathbb ER_l^2=\sum_b r_b^2\ge\beta^2$.  Since $R_l\le M$,
\[
 \mathbb P(R_l\ge\beta/2)\ge\frac{\beta^2}{2M^2}=h.
\]
On this event $U_l,V_l\ge\beta/4$.  For the first time $U_i<c_*$,
condition on the path up to that time.  The squared supermartingale
inequality gives
\[
 \mathbb P\bigl(U_i<c_*\text{ for some }i, U_l\ge\beta/4\bigr)
 \le\frac{c_*^2}{(\beta/4)^2}.
\]
This follows by summing the conditional estimate over the disjoint
first-time events; the same argument applies to $V$.  Our choice of
$c_*$ makes twice this bound at most $h/2$.  Hence the event
\begin{equation}\label{log:good-event}
 G=\{R_l\ge\beta/2,\ U_i,V_i\ge c_*\text{ for every }i\}
\end{equation}
has probability at least $h/2$.

We have retained a definite set of paths on which all ratios can be
logged.  It remains to prevent their product from becoming too small.
Define the increasing concave function
\[
 g(t)=\begin{cases}
 \log c_*+(t-c_*)/c_*,&0\le t<c_*,\\
 \log t,&c_*\le t\le M.
 \end{cases}
\]
The processes $A_i=g(U_i)-g(0)$ and $B_i=g(V_i)-g(0)$ are
nonnegative supermartingales bounded by $H_*$.  Indeed, conditional
Jensen followed by the monotonicity of $g$ proves their supermartingale
inequalities.  Let
\[
 Z=\sum_{i<l}\bigl(g(W_{i+1}^{(i)})-g(W_i^{(i)})\bigr).
\]
For either process $D=A$ or $D=B$, write
$d_i^D=D_i-\mathbb E_iD_{i+1}\ge0$.  Telescoping expectations gives
$\sum_i\mathbb E d_i^D\le H_*$.  Expanding squared increments gives
\begin{equation}\label{log:quadratic-increments}
 \begin{split}
 \sum_{i<l}\mathbb E(D_{i+1}-D_i)^2
 &=\mathbb E(D_l^2-D_0^2)+2\sum_{i<l}\mathbb E(D_i d_i^D)\\
 &\le3H_*^2.
 \end{split}
\end{equation}
Split each selected increment in $Z$ into its conditional mean and
its centered part.  The absolute expectation of the sum of the means
is at most $2H_*$.  The centered parts are martingale differences,
since the choice between $A$ and $B$ is made before the step.  Their
cross terms have zero expectation, and their total variance is at
most $6H_*^2$ by \eqref{log:quadratic-increments}.  Therefore
\begin{equation}\label{log:increment-bound}
 \mathbb E|Z|\le(2+\sqrt6)H_*.
\end{equation}
Markov's inequality and \eqref{log:gamma} give
$\mathbb P(Z<-K)\le h/4$, so
$\mathbb P(G\cap\{Z\ge-K\})\ge h/4$.  On $G$, the product of
ratios in \eqref{log:product} is $e^Z$.  Restricting its expectation
to this last event proves
$r_o\ge (h/4)(\beta/2)e^{-K}=\gamma(M,\beta)$.
\end{proof}

\begin{lemma}[Finite-head gap for the joined-root sum]\label{log:tail-gap}
Let $x,y$ be finitely supported vectors in $X_J$.  Suppose that $x$
is supported in a finite predecessor-closed set $H$ containing the
joining root, and that $y$ vanishes on $H$.  If
\[
 N_J(x+y),N_J(x-y)\le2,
 \qquad N_J(y)\ge\zeta,\qquad 0<\zeta\le2,
\]
then
\[
 \frac{N_J(x+y)+N_J(x-y)}2
 \ge N_J(x)+2\gamma(2,\zeta/4).
\]
\end{lemma}

\begin{proof}
Write $o=\rho$ for the joining root.  Use $p=P(x)$ and $q=Q(x)$,
and define $r$ by
\[
 p+r=\tfrac12\bigl(P(x+y)+P(x-y)\bigr),\qquad
 q+r=\tfrac12\bigl(Q(x+y)+Q(x-y)\bigr).
\]
The same $r$ occurs because both differences of fields equal $x$.
Convexity gives $r\ge0$.  Take the finite tree consisting of $H$ and
the first vertices outside $H$ on paths to the support of $y$.
Call those exterior vertices $B$.  Preserve the complete subtree
values at them.  The four restricted fields are upper fields; a
base tight inequality remains tight at every internal node because
omitted subtrees have zero base fields.  At the root the endpoint
average gives $p_o+q_o+2r_o\le2$.  At $b\in B$,
\[
 p_b=q_b=0,\qquad r_b=\tfrac12(P_b(y)+Q_b(y)).
\]
On $H$ the two fields of $y$ agree.  Upward induction therefore
bounds both by the Euclidean norm of the boundary sums
$P_b(y)+Q_b(y)$ below the node.  In particular,
\[
 N_J(y)\le2\left(\sum_{b\in B}(P_b(y)+Q_b(y))^2\right)^{1/2}
          =4\left(\sum_{b\in B}r_b^2\right)^{1/2}.
\]
Other leaves, if present, have nonnegative $r_b^2$ and cause no loss.
Lemma~\ref{log:finite-tree} with $M=2$ and $\beta=\zeta/4$ gives
$r_o\ge\gamma(2,\zeta/4)$.  The average norm equals $N_J(x)+2r_o$.
\end{proof}

The next theorem is a separated-family formulation of asymptotic
midpoint uniform convexity; see
\citet[Definition~6.15 and Proposition~6.16]{BLP2026}.
Here the field argument supplies the gap directly, with the explicit
constant from Lemma~\ref{log:finite-tree}.

\begin{theorem}[Separated midpoint families]\label{log:lens-gap}
For every $\varepsilon>0$ there is $\eta(\varepsilon)\in(0,1)$ such
that, for $x,z_1,z_2,\ldots\in X_J$,
\[
 \left.
 \begin{gathered}
 N_J(x+z_j),N_J(x-z_j)\le1\quad(j\ge1),\\
 N_J(z_i-z_j)\ge\varepsilon\quad(i\ne j)
 \end{gathered}
 \right\}
 \quad\Longrightarrow\quad N_J(x)\le1-\eta(\varepsilon).
\]
For $0<\varepsilon\le2$, one can take
$\eta(\varepsilon)=\min\{1/2,\gamma(2,\varepsilon/16)\}$.
\end{theorem}

\begin{proof}
The hypotheses imply $N_J(z_j)\le1$, so for $\varepsilon>2$ there
is no such family and any $\eta\in(0,1)$ works.  Suppose
$0<\varepsilon\le2$ and fix
$0<\delta<\min\{1/3,\varepsilon/8\}$.  Choose a finite vector $x'$
with $N_J(x-x')<\delta$, and let $H$ be the finite predecessor
closure of $\operatorname{supp}x'$ together with the root.  The bounded sequence $\Pi_H z_j$ lies in a finite-dimensional
space.  Hence we can choose $i\ne j$ so that, for
$u=(z_i-z_j)/2$, $N_J(\Pi_Hu)<\delta$.  Then
\[
 N_J(u)\ge\varepsilon/2,
 \qquad N_J(x+u),N_J(x-u)\le1.
\]
For example, $x+u$ is the average of $x+z_i$ and $x-z_j$.
The completed tail $u-\Pi_Hu$ can be approximated, by the density
statement of Section~\ref{rec:norms}, by a finite tail $y$ satisfying
$N_J(y-(u-\Pi_Hu))<\delta$.  Consequently
\[
 N_J(y)\ge\varepsilon/2-2\delta\ge\varepsilon/4,
 \qquad N_J(x'\pm y)\le1+3\delta\le2.
\]
Apply Lemma~\ref{log:tail-gap} with $\zeta=\varepsilon/4$ to obtain
\[
 1+3\delta\ge N_J(x')+2\gamma(2,\varepsilon/16)
 \ge N_J(x)-\delta+2\gamma(2,\varepsilon/16).
\]
Let $\delta\downarrow0$.  This proves the stated conclusion, with
room to choose the displayed $\eta$ strictly below one.
\end{proof}
